\documentclass{article}
\usepackage{amsmath}
\usepackage{amssymb}
\usepackage{amsthm}
\usepackage[T1]{fontenc}


\newtheorem*{conjecture}{Przypuszczenie}
\newtheorem*{theorem}{Twierdzenie}
\newtheorem{lemma}{Lemat}
\renewcommand*{\proofname}{Dowód}


\makeatletter
\newcommand*{\rom}[1]{\expandafter\@slowromancap\romannumeral #1@}
\makeatother


\DeclareMathOperator{\interior}{int}

\begin{document}
\tableofcontents


\section{4.1.1. Definicja pochodnej}
Niech \( f \) będzie funkcją określoną na przedziale otwartym \((a, b)\) o wartościach w \(\mathbb{R}\) i \(x_0 \in (a, b)\). Pochodną funkcji \(f\) w punkcie \(x_0\) nazywamy granicę
\[
\lim_{\Delta x \to 0} \frac{f(x_0 + \Delta x) - f(x_0)}{\Delta x} = \lim_{x \to x_0} \frac{f(x) - f(x_0)}{x - x_0},
\]
o ile ta granica istnieje i jest skończona. Granicę \((1)\) oznaczamy przez \(f'(x_0)\). Jeżeli istnieje pochodna \(f'(x_0)\), to mówimy, że funkcja \(f\) jest \emph{różniczkowalna} w \(x_0\).

Liczbę \(\Delta x\) nazywamy przyrostem argumentu funkcji, a liczbę \(\Delta f = f(x_0 + \Delta x) - f(x_0)\) nazywamy przyrostem wartości funkcji. Liczbę \(df = f'(x_0)\Delta x\) nazywamy \emph{różniczką funkcji} \(f\) w punkcie \(x_0\), odpowiadającą przyrostowi \(\Delta x\).

Niech \(o(t)\) oznacza funkcję określoną w sąsiedztwie punktu \(0\) o wartościach w \(\mathbb{R}\) taką, że
\[
\lim_{t \to 0} \frac{o(t)}{t} = 0.
\]
Wtedy funkcję \(o(t)\) nazywamy \emph{małą rzędu wyższego niż \(t\)}. Na przykład funkcja \(x = t^2\) jest typu \(o(t)\). Funkcja \(f\) ma pochodną \(f'(x_0)\) w punkcie \(x_0\), jeżeli
\[
f(x_0 + \Delta x) - f(x_0) = f'(x_0)\Delta x + o(\Delta x).
\]
Istotnie, ponieważ
\[
f'(x_0) = \lim_{\Delta x \to 0} \frac{f(x_0 + \Delta x) - f(x_0)}{\Delta x},
\]
więc
\[
\lim_{\Delta x \to 0} \frac{f(x_0 + \Delta x) - f(x_0) - f'(x_0)\Delta x}{\Delta x} = 0,
\]
a stąd wynika natychmiast wzór \((3)\). Ponieważ wzór \((3)\) można zapisać w postaci
\[
\Delta f = df(x_0)(\Delta x) + o(\Delta x),
\]
a różniczka \(df(x_0)(\Delta x) = f'(x_0)\Delta x\) zależy liniowo od \(\Delta x\), więc różniczkę możemy traktować jako część liniową przyrostu funkcji.

\newpage
\begin{equation*}
    f'(x_0) = \lim_{\Delta x \to 0} \frac{f(x_0 + \Delta x) - f(x_0)}{\Delta x}
\end{equation*}
Odejmuję \(f'(x_0)\) obustronnie i wychodzę na:

\begin{equation*}
    \Bigg(\lim_{\Delta x \to 0} \frac{f(x_0 + \Delta x) - f(x_0)}{\Delta x}\Bigg) - f'(x_0) = 0
\end{equation*}

Wrzucam \(f'(x_0)\) do granicy (w sumie to nie wiem czemu mogę skoro \(f'(x_0)\) to pod spodem granica z delta x0 dążąca do 0):

\begin{equation*}
    \lim_{\Delta x \to 0} \frac{f(x_0 + \Delta x) - f(x_0)}{\Delta x} - f'(x_0)\frac{\Delta x}{\Delta x} = 0
\end{equation*}

Tzn.

\begin{equation*}
    \lim_{\Delta x \to 0} \frac{f(x_0 + \Delta x) - f(x_0) - f'(x_0)\Delta x}{\Delta x} = 0
\end{equation*}

Teraz zapisuję \(0\) jako:

\begin{equation*}
    \lim_{\Delta x \to 0} \frac{f(x_0 + \Delta x) - f(x_0) - f'(x_0)\Delta x}{\Delta x} = 0 = \lim_{\Delta x \to 0} \frac{o(\Delta x)}{\Delta x}
\end{equation*}

I teraz nie wiem co się dzieje, ale to mi wygląda na OBUSTRONNE zdjęcie nagle znaku granicy i zostawieniu pod spodem wyrażeń a potem mnożenie obustronne przez \(\Delta x\):

\begin{equation*}
    \frac{f(x_0 + \Delta x) - f(x_0) - f'(x_0)\Delta x}{\Delta x} = \frac{o(\Delta x)}{\Delta x} 
\end{equation*}
\begin{equation*}
    f(x_0 + \Delta x) - f(x_0) - f'(x_0)\Delta x = o(\Delta x)
\end{equation*}
\begin{equation*}
    f(x_0 + \Delta x) - f(x_0) = o(\Delta x) + f'(x_0)\Delta x
\end{equation*}

\newpage
\begin{theorem}
    Niech \(f: (a, b) \rightarrow \mathbb{R}\) i \(x_0 \in (a, b)\). Funkcja \(f\) ma pochodną \(f'(x_0)\) w punkcie \(x_0\) wtedy i tylko wtedy
    gdy istnieje liczba $A$ taka, że funkcja \(R(\Delta x) := f(x_0 + \Delta x) - f(x_0) - A(\Delta x)\) jest \(o(\Delta x)\).
\end{theorem}

\begin{proof}
    Zakładamy, że \(f\) ma pochodną \(f'(x_0)\) w punkcie \(x_0\). 
    Niech \(R(\Delta x) := f(x_0 + \Delta x) - f(x_0) - A(\Delta x)\). Pokażemy, że funkcja ta jest typu \(o(\Delta x)\) dla pewnej liczby \(A\).
    Zapisujemy:
    \begin{equation*}
        \lim_{\Delta x \to 0} \frac{f(x_0 + \Delta x) - f(x_0) - A \Delta x}{\Delta x} = \lim_{\Delta x \to 0} \frac{f(x_0 + \Delta x) - f(x_0)}{\Delta x} - \lim_{\Delta x \to 0} \frac{A \Delta x}{\Delta x}
    \end{equation*}

    Z definicji pochodnej:

    \begin{equation*}
        = f'(x_{0}) - \lim_{\Delta x \to 0} \frac{A \Delta x}{\Delta x} = f'(x_{0}) - A 
    \end{equation*}
    I aby \(R(\Delta x)\) była \(o(\Delta x)\) potrzeba i wystarcza aby \(A = f'(x_{0})\).
    \newline
    W drugą stronę, załóżmy, że \(R(\Delta x) := f(x_0 + \Delta x) - f(x_0) - A \Delta x \) jest typu \(o(\Delta x)\) dla pewnej liczby \(A\). Pokażemy, że pochodna \(f'(x_{0})\) istnieje i jest równa \(A\).


    Mamy zatem:

    \begin{equation*}
        f(x_0 + \Delta x) - f(x_0)  = R(\Delta x)
    \end{equation*}
    \begin{equation*}
        f(x_0 + \Delta x) - f(x_0) = R(\Delta x) + A \Delta x
    \end{equation*}

    Dzielimy obustronnie przez \(\Delta x\) i przechodzimy do granicy w zerze:

    \begin{equation*}
        \lim_{\Delta x \to 0} \frac{f(x_0 + \Delta x) - f(x_0)}{\Delta x} = \lim_{\Delta x \to 0} \frac{R(\Delta x)}{\Delta x} +  \lim_{\Delta x \to 0} A
    \end{equation*}

    Funkcja \(R(\Delta x)\) jest typu \(o(\Delta x)\) zatem:

    \begin{equation*}
        \lim_{\Delta x \to 0} \frac{R(\Delta x)}{\Delta x} +  \lim_{\Delta x \to 0} A = 0 + \lim_{\Delta x \to 0} A = 0 + A = A
    \end{equation*}

    I z definicji pochodnej:

    \begin{equation*}
        f'(x_{0}) = \lim_{\Delta x \to 0} \frac{f(x_0 + \Delta x) - f(x_0)}{\Delta x} = A
    \end{equation*}
\end{proof}

\end{document}